\section{What We Must Learn to Release}

St.~Benedict commands daily remembrance of death among the tools of good works in chapter~4 of his Rule. Around it stand fear of judgment, desire for eternal life, vigilance over one's conduct, prayer, confession of faults, reconciliation, and the command never to despair of God's mercy. Elsewhere in that same chapter are care for the poor, visiting the sick, burial of the dead, and consolation of the sorrowful. The monk who remembers his own death is also commanded to attend to the suffering and death of others.

Benedict ends his list by forbidding despair of God's mercy. The reader who examines his conduct is therefore also commanded to hope. An unacknowledged sin still needs to be confessed, and death may come before the occasion we have imagined for doing so. We can remember that danger and seek forgiveness without concluding that our repeated failures have put us beyond mercy. Benedict places remembrance of death among practices to be undertaken daily, as part of a life in which repentance continues.

One of the hardest attachments to release is our preferred account of ourselves. We can admit general imperfection while defending every particular instance. We excuse impatience because the other person is difficult, concealment because honesty would threaten our position, or disobedience because we believe we know better. These explanations can make faults familiar without making us willing to amend them. In confession we must name what we have done, accepting that the truth may disturb the account we have grown used to giving. A confessor's questions can help us recognize what our excuses have concealed.

Another attachment is the claim to possess the outcome of good work. A task may really have been entrusted to us, and we may have served it faithfully. Yet its future belongs to providence, which includes other workers with different gifts. A parent cannot control the entire life of a child; a teacher cannot complete every student's education; a pastor cannot see the full harvest of his ministry. To prepare for death is partly to learn the limits of stewardship. We labor with care, make responsible provision, and surrender the claim that love requires us to govern everything forever.

Forgiveness has a similar relation to relinquishment. Some injuries continue to demand recognition, protection, or the repair of justice. Forgiving does not make those claims unreal. It renounces the desire to make another person's suffering the final settlement of our own. Where reconciliation is possible, charity moves toward it truthfully. Where contact would renew harm or where the other person refuses peace, the will can still surrender revenge and seek that person's conversion. A dying person should not be forced into an unsafe encounter in order to satisfy someone else's picture of a perfect farewell.

Possessions also teach us where our hearts have settled. Christ joins treasure and heart (Mt 6:19--21). What we cannot willingly share may reveal a deeper dependence than the amount itself suggests. Almsgiving is an act of worshipful freedom: something under our control becomes available for another's need. Responsible care for dependents remains a duty, and generosity must be ordered by justice. But the duty of provision should not become an indefinite excuse for refusing every gift. The distinction becomes clearer when we remember that death will eventually remove even what we now decline to share.

There is also the attachment to being understood. We may leave a conversation unresolved, our intentions misread, or a work imperfectly appreciated. Charity permits us to explain what should be explained, but it cannot promise that every person will finally see us as we wish. Christ knows the truth. Learning to rest in his knowledge loosens the need to compel agreement from every witness. This can make the last years gentler: there are things we may clarify, things we must repair, and things we may leave with God.

These relinquishments are learned through repeated acts, and repetition should not discourage us. A single generous decision does not abolish possessiveness, nor does one patient response remove the desire to control. Grace forms habits through the return to what is good. Each renewed act places a little more of life under charity. When age or illness narrows the field of action, the same charity can continue within a smaller field. A grateful word, a patient silence, or the acceptance of needed help may then become the full work of a day.

